% !TeX program = tectonic
\documentclass[11pt,a4paper]{article}
\usepackage{fontspec}
\usepackage[russian]{babel}
\babelfont{rm}[Language=Default]{Liberation Serif}
\babelfont[Language=Default]{sf}{Carlito}
\babelfont[Language=Default]{tt}{DejaVu Sans Mono}
\usepackage{amsmath,amssymb,amsthm}
\usepackage{geometry}
\geometry{a4paper, top=2.4cm, bottom=2.4cm, left=2.4cm, right=2.4cm}
\usepackage{booktabs,tabularx,enumitem,xcolor,tcolorbox}
\tcbuselibrary{breakable,skins}
\definecolor{linkblue}{HTML}{1F4E79}
\definecolor{statgreen}{HTML}{2F6B3A}
\definecolor{goldacc}{HTML}{8B7E5A}
\usepackage[colorlinks=true,linkcolor=linkblue,urlcolor=linkblue,unicode=true]{hyperref}
\theoremstyle{plain}
\newtheorem{theorem}{Теорема}
\newtheorem{lemma}[theorem]{Лемма}
\newtheorem{proposition}[theorem]{Предложение}
\newtheorem{corollary}[theorem]{Следствие}
\theoremstyle{definition}
\newtheorem{definition}[theorem]{Определение}
\theoremstyle{remark}
\newtheorem{remark}[theorem]{Замечание}
\newtcolorbox{statusbox}[1][]{colback=statgreen!5,colframe=statgreen!75,
  title={\textbf{Сводка доказанного:} #1},fonttitle=\small,boxrule=0.7pt,arc=2pt,breakable}
\newtcolorbox{databox}[1][]{colback=goldacc!6,colframe=goldacc!80,
  title={\textbf{Данные протокола:} #1},fonttitle=\small,boxrule=0.7pt,arc=2pt,breakable}
\newtcolorbox{verifybox}[1][]{colback=linkblue!5,colframe=linkblue!80,
  title={\textbf{Верификация:} #1},fonttitle=\small,boxrule=0.7pt,arc=2pt,breakable}
\widowpenalty=10000
\clubpenalty=10000
\setlength{\parskip}{0.35em}
\newcommand{\CC}{\mathbb{C}}
\newcommand{\RR}{\mathbb{R}}
\newcommand{\ZZ}{\mathbb{Z}}
\newcommand{\QQ}{\mathbb{Q}}
\newcommand{\Ga}{\Gamma}
\newcommand{\Bfun}{\mathrm{B}}
\newcommand{\im}{\operatorname{Im}}
\newcommand{\re}{\operatorname{Re}}
\newcommand{\lcm}{\operatorname{lcm}}
\newcommand{\SNF}{\operatorname{SNF}}
\newcommand{\Jac}{\operatorname{Jac}}
\newcommand{\rank}{\operatorname{rank}}
\newcommand{\Ind}{\operatorname{Index}}
\newcommand{\disc}{\operatorname{disc}}
\begin{document}
\begin{center}
{\LARGE\bfseries Сертификат E: завершимость потока}\\[0.4em]
{\large точное время t* = lcm и терминальный цикл}\\[0.6em]
{\normalsize Исаев Исхак Хамзатович}\\[0.2em]
{\small Программа <<Динамический принцип>> --- лаборатория \texttt{hodge-laboratory}}\\[0.15em]
{\small Отдельная монография теоремы · Монография 13/16}
\end{center}
\vspace{0.4em}
{\color{linkblue}\hrule height 0.8pt}
\vspace{0.8em}

\section{Постановка}

Сертификат E доказывает завершимость дискретного потока конвейера и
определяет точное время завершения. Любой эксперимент лаборатории
заведомо обрывается --- фундаментальное свойство устойчивости.

\begin{theorem}[E1--E4: завершимость]\label{th:Estand}
\begin{enumerate}[leftmargin=2em]
\item[(E1)] Пространство состояний потока конечно.
\item[(E2)] Существует моновариант открытия: каждая итерация либо
открывает новое состояние, либо переводит поток в терминальный
режим.
\item[(E3)] По принципу Дирихле поток завершается; терминальное
состояние --- цикл.
\item[(E4)] Точное время завершимости
$t^*=\lcm\bigl(W/\gcd(a,W),\,H/\gcd(b,H)\bigr)$.
\end{enumerate}
\end{theorem}

\begin{proof}
(E1) Состояния --- пары (позиция, фаза) на конечной решётке с
конечной группой фаз: не более $4WH$ состояний.
(E2) Моновариант --- множество посещённых ячеек: строго растёт в
режиме открытия, ограничен размером сетки.
(E3) Строгий рост конечного моноварианта обрывается; обрыв ---
циклическое повторение конфигурации по определению перехода.
(E4) Возврат по горизонтали --- через $W/\gcd(a,W)$ шагов, по
вертикали --- через $H/\gcd(b,H)$; совместный возврат --- НОК.
Подстановка эталона ($48\times48$, шаг $(1,1)$) даёт $t^*=48$ ---
и поток завершается ровно за $48$ шагов с полной сверкой
координатной кривой $48/212/432/114$.
\end{proof}

\begin{databox}[восемь пунктов протокола]
E-1: перебор $64$ пар шага --- моновариант растёт, поток
завершается. E-2: конвейер $48\times48$: $t^*=48$; сверка
$48/212/432/114$ полная. E-3: поправочный шаг $(7,22)$ ---
завершимость на правилах поправок. E-4: терминальный цикл длины
$4$ --- $\mu_4$-орбита шага. E-5: масштабы $48\to384$ ---
формула $t^*$ сохраняется. E-6: K3-слой (отражения с торможением)
завершим. E-7: Клейн-слой: Зингер-порождение порядка $7$. E-8:
Клейн-слой: $\mu_4$-орбиты длины $4$. Вердикт: $8/8$ PASS.
\end{databox}

\begin{remark}[цикл --- теорема, не гипотеза]
Циклическая структура терминального состояния --- следствие
конечности и периодичности шага; длина цикла $4$ --- в точности
$\mu_4$-орбита. Это пример теоремы о \emph{структуре} терминального
режима, а не только о его существовании.
\end{remark}

\begin{statusbox}[итог]
Сертификат E принят ($8/8$): поток завершается, время завершимости
вычисляется точно, терминальная структура описана.
\end{statusbox}

\begin{verifybox}[как воспроизвести]
\texttt{python3 laboratory.py} $\to$ <<Сертификаты $\to$ E>>;
конструктор: произвольная сетка $W\times H$ и шаг $(a,b)$ ---
лаборатория печатает $t^*$ и проверяет перебором.
\end{verifybox}

\vfill
{\color{linkblue}\hrule height 0.6pt}
\vspace{0.5em}
{\small\color{gray}
Лицензия: индивидуальная исключительная лицензия Исаева Исхака Хамзатовича.
Все права на вычисления, формулы и тексты принадлежат автору.
Верификация: \texttt{python3 laboratory.py} (пункт <<Стенды>>/<<Сертификаты>>),
Lean: \texttt{verification/lean/HodgeLaboratory.lean}.
}
\end{document}
